\begin{codebox}
\codeline{\textcolor{cmtgray}{\#\ 一阶逻辑表达式示例}}
\codeline{\textcolor{cmtgray}{\#\ First-Order\ Logic\ Expressions\ in\ Ontology\ Engineering}}
\codeline{\textcolor{cmtgray}{\#}}
\codeline{\textcolor{cmtgray}{\#\ 一阶逻辑（First-Order\ Logic,\ FOL）是本体表达的重要基础}}
\codeblank{}
\codeline{\textcolor{cmtgray}{\#\ ==============================}}
\codeline{\textcolor{cmtgray}{\#\ 量词与逻辑联结词}}
\codeline{\textcolor{cmtgray}{\#\ ==============================}}
\codeblank{}
\codeline{\textcolor{cmtgray}{\#\ 量词（Quantifiers）：}}
\codeline{\textcolor{cmtgray}{\#\ \ \ 全称量词\ \(\forall\)（对于所有）}}
\codeline{\textcolor{cmtgray}{\#\ \ \ 存在量词\ \(\exists\)（存在）}}
\codeblank{}
\codeline{\textcolor{cmtgray}{\#\ 逻辑联结词：}}
\codeline{\textcolor{cmtgray}{\#\ \ \ 与\ (\(\land\))、或\ (\(\lor\))、非\ (\(\lnot\))、蕴含\ (\(\rightarrow\))}}
\codeblank{}
\codeline{\textcolor{cmtgray}{\#\ ==============================}}
\codeline{\textcolor{cmtgray}{\#\ 工程约束的一阶逻辑表达}}
\codeline{\textcolor{cmtgray}{\#\ ==============================}}
\codeblank{}
\codeline{\textcolor{cmtgray}{\#\ 规则：所有设备都需要定期维护}}
\codeline{\(\forall\)x:\ Equipment(x)\ \(\rightarrow\)\ \(\exists\)m:\ Maintenance(m)\ \(\land\)\ requires(x,\ m)\ \(\land\)\ isPeriodic(m)}
\codeblank{}
\codeline{\textcolor{cmtgray}{\#\ 规则：容量大的故障设备优先维修}}
\codeline{\(\forall\)a,b:\ (Equipment(a)\ \(\land\)\ Equipment(b)\ \(\land\)\ capacity(a)\ >\ capacity(b)\ \(\land\)\ hasFault(a))}
\codeline{\hspace*{4\ttcw}\(\rightarrow\)\ priority(repair(a))\ >\ priority(repair(b))}
\codeblank{}
\codeline{\textcolor{cmtgray}{\#\ 规则：同一时刻设备只能执行一个任务}}
\codeline{\(\forall\)\ e,\ t,\ task1,\ task2:\ (executes(e,\ task1,\ t)\ AND\ executes(e,\ task2,\ t))}
\codeline{\hspace*{4\ttcw}\(\rightarrow\)\ (task1\ =\ task2)}
\codeblank{}
\codeline{\textcolor{cmtgray}{\#\ 规则：热处理必须在焊接之后进行}}
\codeline{\(\forall\)\ p1,\ p2:\ (p1.type\ =\ "HeatTreatment"\ AND\ p2.type\ =\ "Welding")}
\codeline{\hspace*{4\ttcw}\(\rightarrow\)\ precedes(p2,\ p1)}
\codeblank{}
\codeline{\textcolor{cmtgray}{\#\ ==============================}}
\codeline{\textcolor{cmtgray}{\#\ 非单调推理\ (Non-monotonic\ Reasoning)}}
\codeline{\textcolor{cmtgray}{\#\ ==============================}}
\codeline{\textcolor{cmtgray}{\#\ 默认规则：设备一般是可用的，除非明确知道其发生故障}}
\codeblank{}
\codeline{Equipment(x)\ :\ available(x)\ /\ available(x)}
\codeline{\textcolor{cmtgray}{\#\ 规则解读：如果x是设备，且"x可用"这一假设是一致的}}
\codeline{\textcolor{cmtgray}{\#\ （即知识库中没有x故障的信息），则默认假设x可用。}}
\codeblank{}
\codeline{\textcolor{cmtgray}{\#\ 当知识库中添加\ hasFault(Lathe\_001)\ 后：}}
\codeline{\textcolor{cmtgray}{\#\ \(\rightarrow\)\ "Lathe\_001可用"\ 的假设不再成立}}
\codeline{\textcolor{cmtgray}{\#\ \(\rightarrow\)\ 之前的推论被推翻（撤销推理）}}
\codeline{\textcolor{cmtgray}{\#\ \(\rightarrow\)\ 对于没有故障记录的设备，仍然默认其可用}}
\end{codebox}
